\import{set.tex}
\import{nat.tex}

\section{Propositional formulas}


\begin{datatype}\label{propform}
    %! infixr 0
    Define $\propform$ inductively as follows.
    \begin{enumerate}
        \item $\propbot \in \propform$.
        \item $\propvar{n} \in \propform$ for $n \in \naturals$.
        \item $(p \propto q) \in \propform$ for $p \in \propform$ and $q \in \propform$.
    \end{enumerate}
\end{datatype}

\begin{abbreviation}\label{propnot}
    $\propnot{p} = (p \propto \propbot)$.
\end{abbreviation}

\begin{proposition}\label{propform_bot_intro}
    $\propbot \in \propform$.
\end{proposition}
\begin{proof}
    Follows by \obv{propform_propbot_intro}.
\end{proof}

\begin{proposition}\label{propform_var_intro}
    Suppose $n \in \naturals$.
    Then $\propvar{n} \in \propform$.
\end{proposition}
\begin{proof}
    Follows by \obv{propform_propvar_intro}.
\end{proof}

\begin{proposition}\label{propform_imp_intro}
    Suppose $p \in \propform$ and $q \in \propform$.
    Then $(p \propto q) \in \propform$.
\end{proposition}
\begin{proof}
    Follows by \obv{propform_propto_intro}.
\end{proof}

\begin{proposition}\label{propnot_intro}
    Suppose $p \in \propform$.
    Then $\propnot{p} \in \propform$.
\end{proposition}
\begin{proof}
    Follows by \obv{propform_propbot_intro,propform_propto_intro}.
\end{proof}

Predicate-level inductive derivability is deferred until Felix supports
inductively defined predicates. For now, the generated datatype induction
principle is available for reasoning about formulas themselves.

\begin{proposition}\label{propform_induction_sanity}
    For every $x \in \propform$ we have $x \in \propform$.
\end{proposition}
\begin{proof}
    Follows by \obv{propform_induct,propform_propbot_intro,propform_propvar_intro,propform_propto_intro}.
\end{proof}
